% =============================================================================
%  Togo_Sophos_DeclarationPI.tex
%  Équipe Sophos, solution SOLIDA
%  Hackathon National CIF / DigiCoop-WA+, Thématique 02, Scoring Microcrédit (Togo)
%
%  DÉCLARATION DE PROPRIÉTÉ INTELLECTUELLE, article 10 de la Charte CIF.
%  Pièce à signer par chaque membre de l'équipe et à joindre aux livrables.
%  Elle déclare les éléments créés par l'équipe et la licence libre accordée
%  (Apache-2.0), ainsi que les composants tiers et logiciels libres intégrés à la
%  solution, avec leur licence exacte, et garantit leur compatibilité avec la
%  cession prévue par la Charte.
%
%  Sources : Charte de propriété intellectuelle CIF signée ; réponses officielles
%  au suivi des questions candidats ; fiche de présentation de l'équipe ; état
%  réel du dépôt SOLIDA au 2026-09-02. Licences des composants tiers vérifiées à
%  la source (registre npm, PyPI) le 2026-09-02.
%
%  Compilation : pdflatex Togo_Sophos_DeclarationPI.tex  (2 passes)
%  Style : aucun tiret cadratin ; nommage métier en français, technique en anglais.
%  À figer avant signature : étiquette et empreinte du commit Git (section 1.1),
%  date de signature.
% =============================================================================

\documentclass[11pt,a4paper]{article}

\usepackage[utf8]{inputenc}
\usepackage[T1]{fontenc}
\usepackage[french]{babel}
\usepackage[a4paper,top=1.8cm,bottom=1.6cm,left=2.0cm,right=2.0cm]{geometry}
\usepackage{tikz}
\usepackage{tcolorbox}
\usepackage{booktabs}
\usepackage{enumitem}
\setlist[itemize]{label=\textbullet}
\usepackage{fancyhdr}
\usepackage{array}
\usepackage{longtable}
\usepackage{xcolor}
\usepackage[expansion=false]{microtype}

\tcbuselibrary{skins,breakable}

% ---- Palette (identité des documents administratifs de l'équipe) ----
\definecolor{navy}{HTML}{0B1F3A}
\definecolor{gold}{HTML}{C9A961}
\definecolor{lightgold}{HTML}{F5EFE0}
\definecolor{softgray}{HTML}{4A4A4A}

% ---- En-tetes et pieds ----
\pagestyle{fancy}
\fancyhf{}
\renewcommand{\headrulewidth}{0.4pt}
\renewcommand{\footrulewidth}{0.4pt}
\fancyhead[L]{\footnotesize\color{navy}\textsc{Sophos} : Déclaration PI (article 10)}
\fancyhead[R]{\footnotesize\color{navy}Hackathon CIF DigiCoop-WA+ (Togo)}
\fancyfoot[C]{\footnotesize\color{softgray}\thepage}

\newcommand{\hr}{\vspace{2pt}\textcolor{gold}{\rule{\textwidth}{0.4pt}}\vspace{2pt}}

\linespread{0.99}\selectfont
\begin{document}

% ============ BANDEAU DE TITRE ============
\begin{center}
\begin{tikzpicture}
\node[rectangle, rounded corners=3pt, fill=navy, text=white, minimum width=\textwidth,
      inner sep=12pt, align=center] (box) {
  \begin{minipage}{0.93\textwidth}
    \centering
    {\Large\bfseries Déclaration de propriété intellectuelle}\\[4pt]
    {\normalsize\color{gold}\bfseries Article 10 de la Charte de propriété intellectuelle CIF}\\[4pt]
    {\footnotesize Équipe Sophos $\cdot$ Solution SOLIDA $\cdot$ Thématique 02, Scoring Microcrédit}\\[2pt]
    {\footnotesize Hackathon National d'Innovation CIF / DigiCoop-WA+, Lomé, 13 au 15 septembre 2026}
  \end{minipage}
};
\end{tikzpicture}
\end{center}

\vspace{0.6em}

\section*{Objet}

La présente déclaration répond à l'article 10 de la Charte de propriété intellectuelle de la CIF.
Elle identifie les éléments créés par l'équipe Sophos et la licence libre que l'équipe leur
applique, puis les composants tiers et logiciels libres intégrés à la solution SOLIDA, avec leur
licence exacte, et garantit leur compatibilité avec la cession et l'exploitation prévues par la
Charte. Elle est signée par chaque membre de l'équipe et jointe aux livrables.

\section*{1. Éléments préexistants de l'équipe et licence accordée}

\subsection*{1.1 Périmètre concerné}

Le présent article porte sur le \textbf{périmètre préexistant} : l'état du dépôt Git de la solution
SOLIDA figé par l'étiquette \texttt{socle-pre-hackathon}, empreinte de commit
\texttt{ccad4d7a219dba59ccd2353684d73f024cc7641f}. Les développements réalisés pendant le
hackathon n'en font pas partie (voir 1.4).

\subsection*{1.2 Contenu du périmètre préexistant}

Éléments conçus et développés par l'équipe Sophos avant l'épreuve des 72 heures, présents dans le
dépôt. Ils constituent sa création propre au sens de l'article 9 de la Charte.

\begin{center}
\footnotesize
\renewcommand{\arraystretch}{1.25}
\begin{tabular}{@{}p{3.9cm} p{8.6cm} p{2.3cm}@{}}
\toprule
\textbf{Élément} & \textbf{Description} & \textbf{Licence} \\
\midrule
Interface agent SOLIDA & Application web Next.js : connexion, recherche, dossier du sociétaire, nouvelle demande, résultat de scoring, fiche de justification, registre des décisions, paramétrage de la grille & Apache-2.0 \\
Backend applicatif SOLIDA & API FastAPI : lecture du simulateur, scoring, moteur de décision, authentification, registre, génération et archivage de la fiche PDF & Apache-2.0 \\
Simulateur de données synthétiques & Générateur Python de données fictives d'une coopérative et chargement PostgreSQL & Apache-2.0 \\
\bottomrule
\end{tabular}
\end{center}

\subsection*{1.3 Licence accordée}

L'équipe Sophos place ce périmètre préexistant, code source et documentation qui l'accompagne,
sous \textbf{Apache License 2.0} (fichier \texttt{LICENSE} du dépôt). Cette licence accorde à la
CIF, comme à tout tiers, le droit mondial, gratuit et pérenne d'utiliser, reproduire, modifier,
distribuer et exploiter ces éléments, y compris commercialement, avec la concession de brevet
qu'elle prévoit. Conformément à l'article 10 de la Charte, la CIF bénéficie en outre d'une licence
non exclusive, mondiale, transférable et gratuite, pour toute la durée de protection, sur ces
mêmes éléments dans le cadre de la Solution. Les fichiers \texttt{LICENSE} et \texttt{NOTICE} sont
présents dans le dépôt. En signant la présente déclaration, chaque membre de l'équipe confirme
cette concession de licence.

\textbf{Nom \og SOLIDA \fg{}.} La dénomination et la charte graphique ont été retenues par
l'équipe pour identifier la solution. L'équipe ne revendique aucune marque déposée et en concède
l'usage libre et gratuit à la CIF. L'article 5.4 de la Charte prévoyant que les signes créés pour
identifier la Solution sont cédés à la CIF, l'équipe ne s'y oppose pas.

\subsection*{1.4 Articulation avec les développements du hackathon}

Les développements réalisés pendant l'épreuve des 72 heures et pendant le suivi post-hackathon du
lauréat \textbf{appartiennent définitivement à la CIF}, par cession exclusive et définitive
(articles 4 à 6 de la Charte). L'équipe n'en conserve aucun droit d'exploitation autonome
(article 14). La licence libre décrite au 1.3 ne couvre \textbf{que} le périmètre préexistant
défini au 1.1 ; elle ne s'étend pas à ces développements et n'en affecte pas la cession.

\section*{2. Composants tiers et logiciels libres intégrés}

Dépendances de \textbf{premier niveau}, avec version épinglée et licence vérifiée à la source le
2026-09-02. L'inventaire complet des dépendances directes et transitives figure au dépôt :
\texttt{THIRD\_PARTY\_NOTICES.md}, \texttt{docs/livrables/sbom-frontend.txt} (855 paquets),
\texttt{docs/livrables/sbom-backend.txt} (81 paquets), \texttt{frontend/package-lock.json} et
\texttt{backend/uv.lock}.

{\footnotesize
\renewcommand{\arraystretch}{1.12}
\begin{longtable}{@{}p{4.7cm} p{2.0cm} p{3.1cm} p{5.5cm}@{}}
\toprule
\textbf{Composant} & \textbf{Version} & \textbf{Licence} & \textbf{Usage dans la solution} \\
\midrule
\endfirsthead
\toprule
\textbf{Composant} & \textbf{Version} & \textbf{Licence} & \textbf{Usage dans la solution} \\
\midrule
\endhead
\bottomrule
\endfoot
\bottomrule
\endlastfoot
\multicolumn{4}{@{}l}{\textbf{Frontend, exécution}} \\
\addlinespace[2pt]
next & 16.3.0 & MIT & Framework React, rendu serveur et routage \\
react & 19.2.8 & MIT & Composants d'interface \\
react-dom & 19.2.8 & MIT & Rendu DOM de React \\
radix-ui & 1.6.7 & MIT & Primitives d'interface accessibles \\
recharts & 3.10.1 & MIT & Graphiques d'épargne et de simulation \\
framer-motion & 12.43.0 & MIT & Animations d'interface \\
lucide-react & 1.28.0 & ISC & Jeu d'icônes \\
cmdk & 1.1.1 & MIT & Palette de commandes et recherche \\
sonner & 2.0.7 & MIT & Notifications transitoires \\
clsx & 2.1.1 & MIT & Concaténation conditionnelle de classes CSS \\
tailwind-merge & 3.6.0 & MIT & Fusion des classes utilitaires \\
class-variance-authority & 0.7.1 & Apache-2.0 & Composition des variantes de style \\
\addlinespace[3pt]
\multicolumn{4}{@{}l}{\textbf{Frontend, outillage de développement et de test (hors exécution)}} \\
\addlinespace[2pt]
tailwindcss & 4.3.3 & MIT & Feuilles de style utilitaires \\
@tailwindcss/postcss & 4.3.3 & MIT & Intégration PostCSS \\
tw-animate-css & 1.4.0 & MIT & Utilitaires d'animation \\
typescript & 5.9.3 & Apache-2.0 & Typage statique \\
eslint & 9.39.5 & MIT & Analyse statique \\
eslint-config-next & 16.3.0 & MIT & Règles ESLint Next.js \\
prettier & 3.9.6 & MIT & Formatage du code \\
husky & 9.1.7 & MIT & Hooks Git \\
lint-staged & 17.3.0 & MIT & Vérifications de pré-commit \\
shadcn & 4.16.1 & MIT & CLI de génération de composants (copiés dans le dépôt sous MIT) \\
vitest & 4.1.10 & MIT & Tests unitaires \\
@types/node & 22.20.1 & MIT & Déclarations de types \\
@types/react & 19.2.18 & MIT & Déclarations de types \\
@types/react-dom & 19.2.4 & MIT & Déclarations de types \\
\addlinespace[3pt]
\multicolumn{4}{@{}l}{\textbf{Backend et API, exécution}} \\
\addlinespace[2pt]
fastapi & 0.141.1 & MIT & Framework de l'API de scoring \\
pydantic & 2.13.4 & MIT & Validation des schémas de données \\
pydantic-settings & 2.14.2 & MIT & Chargement de la configuration \\
sqlalchemy & 2.0.51 & MIT & Accès relationnel (ORM et Core) \\
alembic & 1.18.5 & MIT & Migrations de la base SOLIDA \\
psycopg[binary] & 3.3.4 & LGPL-3.0-only & Pilote PostgreSQL (voir section 3) \\
fastapi-users[sqlalchemy] & 15.0.5 & MIT & Authentification et comptes agents \\
argon2-cffi & 25.1.0 & MIT & Hachage des mots de passe (Argon2) \\
uvicorn[standard] & 0.52.1 & BSD-3-Clause & Serveur ASGI \\
weasyprint & 69.0 & BSD-3-Clause & Rendu PDF de la fiche de justification \\
jinja2 & 3.1.6 & BSD-3-Clause & Gabarit HTML de la fiche \\
minio & 7.2.20 & Apache-2.0 & Client S3 des fiches archivées \\
\addlinespace[3pt]
\multicolumn{4}{@{}l}{\textbf{Backend, outillage de développement et de test (hors exécution)}} \\
\addlinespace[2pt]
ruff & 0.16.1 & MIT & Lint \\
mypy & 2.3.0 & MIT & Vérification de types \\
pytest & 9.1.1 & MIT & Tests \\
pytest-cov & 7.1.0 & MIT & Couverture de tests \\
httpx & 0.28.1 & BSD-3-Clause & Client HTTP de test \\
import-linter & 2.13 & BSD-2-Clause & Contrôle des couches (Clean Architecture) \\
hatchling & 1.x & MIT & Construction du paquet \\
\addlinespace[3pt]
\multicolumn{4}{@{}l}{\textbf{Simulateur de données synthétiques, exécution}} \\
\addlinespace[2pt]
numpy & 2.5.1 & BSD-3-Clause & Calcul numérique \\
pandas & 3.0.5 & BSD-3-Clause & Manipulation des jeux de données \\
pyyaml & 6.0.3 & MIT & Lecture de la configuration du générateur \\
sqlalchemy & 2.0.51 & MIT & Chargement vers PostgreSQL \\
psycopg2-binary & 2.9.12 & LGPL-3.0-or-later & Pilote PostgreSQL du chargement (voir section 3) \\
pyarrow & 25.0.0 & Apache-2.0 & Lecture et écriture au format Parquet \\
\addlinespace[3pt]
\multicolumn{4}{@{}l}{\textbf{Infrastructure, images Docker de base}} \\
\addlinespace[2pt]
node:22-slim & 22 & MIT (Node.js), base Debian & Image d'exécution du frontend \\
python:3.12-slim & 3.12 & Python-2.0 (PSF), base Debian & Image d'exécution du simulateur \\
ghcr.io/astral-sh/uv & python3.12-trixie-slim & MIT OR Apache-2.0, base Debian & Image du backend \\
postgres & 16.14 & PostgreSQL License, base Debian & Bases SOLIDA et CORE-SIM \\
nginx & 1.30.4-alpine & BSD-2-Clause (nginx), base Alpine & Reverse proxy, point d'entrée réseau \\
chrislusf/seaweedfs & 4.40 & Apache-2.0 & Stockage objet compatible S3 \\
\end{longtable}
}

\noindent \textbf{Noms complets des licences citées} (identifiant SPDX, puis intitulé officiel) :

\begin{center}
\footnotesize
\renewcommand{\arraystretch}{1.15}
\begin{tabular}{@{}p{3.6cm} p{12.2cm}@{}}
\toprule
\textbf{Identifiant SPDX} & \textbf{Intitulé officiel} \\
\midrule
\texttt{MIT} & MIT License \\
\texttt{Apache-2.0} & Apache License 2.0 \\
\texttt{BSD-3-Clause} & BSD 3-Clause \og New \fg{} or \og Revised \fg{} License \\
\texttt{BSD-2-Clause} & BSD 2-Clause \og Simplified \fg{} License \\
\texttt{ISC} & ISC License \\
\texttt{LGPL-3.0-only} & GNU Lesser General Public License version 3.0 only \\
\texttt{LGPL-3.0-or-later} & GNU Lesser General Public License version 3.0 or later \\
\texttt{Python-2.0} & Python Software Foundation License 2.0 \\
\texttt{PostgreSQL} & The PostgreSQL License (permissive, type BSD) \\
\texttt{MIT OR Apache-2.0} & double licence : MIT License ou Apache License 2.0, au choix \\
\bottomrule
\end{tabular}
\end{center}

\noindent \textbf{Dépendances transitives.} Elles relèvent des mêmes familles de licences
permissives : MIT License, BSD 2-Clause, BSD 3-Clause, Apache License 2.0, ISC License, BSD Zero
Clause License (\texttt{0BSD}), zlib License, Python Software Foundation License 2.0. Aucune
dépendance applicative, directe ou transitive, n'est sous copyleft fort (GNU GPL, GNU AGPL).

\noindent \textbf{Images Docker de base.} Ce sont des distributions système standard (Debian,
Alpine), utilisées sans modification. Leur contenu inclut, comme toute distribution Linux, des
paquets sous copyleft dont la GNU GPL ; ils forment l'environnement d'exécution et ne sont ni liés
ni incorporés au code de la Solution.

\section*{3. Note sur la licence LGPL (pilotes PostgreSQL)}

\texttt{psycopg} est sous \textbf{GNU Lesser General Public License version 3.0 seule}
(\texttt{LGPL-3.0-only}) ; \texttt{psycopg2-binary} sous \textbf{GNU LGPL version 3.0 ou
ultérieure} avec exceptions de compilation (\texttt{LGPL-3.0-or-later}). Cette licence
\textbf{accorde} : l'usage commercial et privé sans redevance ; l'intégration dans un logiciel
propriétaire ou fermé \textbf{sans obligation d'en publier le code} ; la modification et la
redistribution de la bibliothèque. Elle \textbf{impose} : en cas de modification des fichiers
internes de la bibliothèque, la publication de ces modifications sous LGPL ; la conservation des
mentions de copyright et de licence. Le logiciel est fourni sans garantie.

Dans SOLIDA, les deux pilotes sont utilisés \textbf{tels quels} (installés par gestionnaire de
paquets, importés sans modification) et liés dynamiquement. Aucune obligation de diffusion du code
de SOLIDA n'en découle. La cession et l'exploitation prévues par la Charte, y compris pour une
solution dont la propriété serait transférée à la CIF, restent possibles : il suffit que le pilote
demeure non modifié et remplaçable, ce que garantit l'accès à la base via SQLAlchemy.

\section*{4. Garantie de compatibilité (articles 9 et 10)}

L'équipe Sophos garantit que la Solution est sa création propre et ne porte pas atteinte aux
droits de tiers. Les travaux ont porté exclusivement sur des données synthétiques ; l'usage
d'assistants de génération de code (Claude Code, Codex) a respecté la Charte, l'originalité de la
création revenant à l'équipe, et n'a transmis aucune donnée personnelle réelle à un service tiers. Les composants tiers listés relèvent de licences
permissives (MIT, BSD 2-Clause, BSD 3-Clause, Apache License 2.0, ISC) ou de copyleft faible
(GNU LGPL 3.0, limité aux pilotes PostgreSQL, employés sans modification), toutes compatibles avec
la cession et l'exploitation prévues par la Charte. Tout composant sous licence incompatible qui
serait identifié avant ou après la remise des livrables sera remplacé et la présente déclaration
mise à jour.

\section*{5. Livrables associés (article 11)}

Seront remis avec la Solution : le code source et objet, les identifiants et accès techniques, les
fichiers \texttt{LICENSE} et \texttt{NOTICE}, le fichier \texttt{THIRD\_PARTY\_NOTICES.md} et les
inventaires de dépendances \texttt{sbom-frontend.txt} et \texttt{sbom-backend.txt} (dossier
\texttt{docs/livrables/}).

\section*{6. Signatures}

Fait à Lomé, le \rule{2.2cm}{0.4pt} septembre 2026.

\smallskip
\noindent Chaque membre reconnaît la présente déclaration et confirme, pour ce qui le concerne, la
concession de licence décrite à la section 1. Signature précédée de la mention \og lu et approuvé \fg{}.

\smallskip
\begin{center}
\footnotesize
\renewcommand{\arraystretch}{1.7}
\begin{tabular}{@{}p{0.6cm} p{4.7cm} p{4.6cm} p{4.0cm}@{}}
\toprule
\textbf{N\textsuperscript{o}} & \textbf{Nom et prénom} & \textbf{Rôle} & \textbf{Signature} \\
\midrule
1 & ADOGLI Jean-Paul & Chef d'équipe, Data Science / Modèle & \\
2 & DJIMEDO Yao-Junior Samuel & Data Engineering & \\
3 & POULI Epiphane-Gédéon & Backend / API et authentification & \\
4 & AKAKPO Charlie & DevOps / Infrastructure et intégration & \\
\bottomrule
\end{tabular}
\end{center}

\begin{center}
{\footnotesize\color{softgray}Équipe Sophos $\cdot$ Déclaration de propriété intellectuelle
(article 10) $\cdot$ Hackathon National CIF / DigiCoop-WA+, Lomé 2026}
\end{center}

\end{document}
